% !TEX program = xelatex
% ============================================================
%  خطة المشتريات — Procurement Plan Template
%  نموذج لتخطيط مشتريات المشروع وتتبع حالة كل بند شراء
%  Compile with: xelatex main.tex (run twice)
% ============================================================
%  Features:
%    - Procurement table: Item | Description | Quantity |
%      Estimated cost | Vendor | Lead time | Status
%    - Total row with cost summary
%    - Notes section
% ============================================================

\documentclass[11pt,a4paper]{article}

\usepackage[a4paper,margin=2.5cm,top=2.5cm,bottom=2.5cm,headheight=16pt]{geometry}
\usepackage{graphicx}
\usepackage{array}
\usepackage{booktabs}
\usepackage{tabularx}
\usepackage{multirow}
\usepackage{enumitem}
\usepackage{float}
\usepackage{titlesec}
\usepackage{fancyhdr}
\usepackage{setspace}
\usepackage{xcolor}
\usepackage{amsmath}

\usepackage{bidi}
\usepackage[no-math]{fontspec}
\usepackage[quiet]{polyglossia}

\setmainfont[Scale=1]{NotoKufiArabic-Regular.ttf}
\newfontfamily\arabicfont[Script=Arabic,Scale=0.95]{NotoKufiArabic-Regular.ttf}
\newfontfamily\englishfont[Scale=0.95]{TeX Gyre Termes}

\setdefaultlanguage[calendar=gregorian,hijricorrection=1]{arabic}
\setotherlanguage[variant=british]{english}

\usepackage[breaklinks,hidelinks]{hyperref}

\addto\captionsarabic{%
  \renewcommand{\contentsname}{الفهرس}%
  \renewcommand{\figurename}{الشكل}%
  \renewcommand{\tablename}{الجدول}%
}

\titleformat{\section}{\Large\bfseries}{\thesection}{0.7em}{}
\titlespacing*{\section}{0pt}{1.4ex plus 0.4ex}{0.9ex plus 0.3ex}

\pagestyle{fancy}
\fancyhf{}
\fancyhead[R]{\small\itshape خطة المشتريات}
\fancyhead[L]{\small اسم المشروع}
\fancyfoot[C]{\small\thepage}
\renewcommand{\headrulewidth}{0.3pt}

\setlist[itemize]{topsep=0.3em, itemsep=0.15em, parsep=0pt}
\setlist[enumerate]{topsep=0.3em, itemsep=0.15em, parsep=0pt}

\newcolumntype{L}[1]{>{\raggedright\arraybackslash}p{#1}}
\newcolumntype{C}[1]{>{\centering\arraybackslash}p{#1}}
\newcolumntype{R}[1]{>{\raggedleft\arraybackslash}p{#1}}

\setstretch{1.3}

\begin{document}

% ============================================================
% TITLE BLOCK
% ============================================================
\begin{center}
{\LARGE\bfseries خطة المشتريات}\\[0.3em]
{\large Procurement Plan}\\[1em]
\end{center}

% ============================================================
% PROJECT INFO
% ============================================================
% TODO: املأ بيانات المشروع.
\begin{table}[H]
\centering
\renewcommand{\arraystretch}{1.4}
\begin{tabularx}{\textwidth}{L{3.5cm} X L{3.5cm} X}
\toprule
\textbf{اسم المشروع:} & --- & \textbf{مدير المشروع:} & --- \\
\textbf{اسم الشركة:} & --- & \textbf{تاريخ التحديث:} & --- \\
\textbf{العملة:} & \LR{USD} & \textbf{إجمالي الميزانية:} & --- \\
\bottomrule
\end{tabularx}
\end{table}

\vspace{1em}

% ============================================================
% 1. PROCUREMENT TABLE
% ============================================================
\section{بنود المشتريات}
\label{sec:procurement}

% TODO: عدّل بنود المشتريات حسب مشروعك.
\begin{table}[H]
\centering
\caption{بنود المشتريات المخططة}
\renewcommand{\arraystretch}{1.4}
\small
\begin{tabularx}{\textwidth}{L{2.8cm} L{3cm} C{1.5cm} R{2.2cm} L{2.5cm} C{1.8cm} L{1.8cm}}
\toprule
\textbf{البند} & \textbf{الوصف} & \textbf{الكمية} & \textbf{التكلفة التقديرية} & \textbf{المورد} & \textbf{مدة التسليم} & \textbf{الحالة} \\
\midrule
أجهزة حاسب & حواسيب محمولة للمطورين & 5 & 7{,}500 & --- & 2 أسبوع & مخطط \\
خوادم & خوادم لإستضافة النظام & 2 & 12{,}000 & --- & 4 أسابيع & قيد الطلب \\
تراخيص برمجيات & تراخيص بيئة التطوير & 10 & 5{,}000 & --- & 1 أسبوع & مكتمل \\
أجهزة شبكة & موجات ومحولات شبكة & 3 & 3{,}000 & --- & 2 أسبوع & مخطط \\
معدات مكتبية & أثاث ومعدات مكتبية & --- & 4{,}000 & --- & 3 أسابيع & مخطط \\
خدمات استشارية & استشارات أمن المعلومات & 1 & 8{,}000 & --- & 6 أسابيع & قيد التقييم \\
مواد استهلاكية & قرطاسية ومواد طباعة & --- & 1{,}500 & --- & 1 أسبوع & مكتمل \\
\midrule
\textbf{الإجمالي} & \multicolumn{2}{l}{} & \textbf{41{,}000} & \multicolumn{3}{l}{} \\
\bottomrule
\end{tabularx}
\end{table}

% ============================================================
% 2. PROCUREMENT SCHEDULE
% ============================================================
\section{جدول المشتريات}
\label{sec:schedule}

% TODO: حدّد الجدول الزمني لكل عملية شراء.
\begin{table}[H]
\centering
\caption{الجدول الزمني للمشتريات}
\renewcommand{\arraystretch}{1.4}
\begin{tabularx}{\textwidth}{L{3cm} L{3cm} L{3cm} X}
\toprule
\textbf{البند} & \textbf{تاريخ الطلب} & \textbf{تاريخ التسليم المتوقع} & \textbf{ملاحظات} \\
\midrule
أجهزة حاسب & --- & --- & طلب عاجل لبدء التطوير \\
خوادم & --- & --- & يتطلب إعداد مسبق لموقع الخوادم \\
تراخيص برمجيات & --- & --- & تم التسليم إلكترونياً \\
أجهزة شبكة & --- & --- & تنسيق مع قسم تقنية المعلومات \\
معدات مكتبية & --- & --- & طلب مجمع لتقليل التكلفة \\
خدمات استشارية & --- & --- & يتطلب تقييم العروض أولاً \\
مواد استهلاكية & --- & --- & شراء دوري حسب الحاجة \\
\bottomrule
\end{tabularx}
\end{table}

% ============================================================
% 3. NOTES
% ============================================================
\section{ملاحظات}
\label{sec:notes}

% TODO: أضف ملاحظات حول خطة المشتريات.
\begin{itemize}
    \item جميع المشتريات تخضع لموافقة مدير المشروع والجهة المالية قبل التنفيذ.
    \item المشتريات التي تتجاوز \LR{10,000 USD} تتطلب عرض سعر تنافسي من ثلاثة موردين على الأقل.
    \item يجب مراعاة مدد التسليم عند جدولة المهام المعتمدة على المشتريات.
    \item يُفضّل التعاقد مع موردين محليين لتقليل مدد التسليم وتكاليف الشحن.
\end{itemize}

\end{document}
